% -------------------------------------------------------------------------

\subsection{Теоретический вопрос II.1: Почему columns of a matrix являются images of basis vectors?}
\textbf{Формулировка вопроса:} Докажите, почему столбцы матрицы линейного отображения в точности равны координатным
представлениям образов базисных векторов.

\textbf{Ответ и строгое доказательство:}

Пусть $T \colon V \to W$ — линейное отображение между векторными пространствами.
Зафиксируем базис $\mathcal{B} = (v_1, \dots, v_n)$ в $V$ и базис $\mathcal{C} = (w_1, \dots, w_m)$ в $W$.
Возьмём вектор $x \in V$.
По теореме о базисе он единственным образом представляется в виде:
\[ x = \sum_{j=1}^n x_j v_j, \quad \text{где } [x]_\mathcal{B} = \begin{pmatrix}
                                                                     x_1 \\ \vdots \\ x_n
\end{pmatrix}. \]
Применим отображение $T$ к вектору $x$.
В силу свойства \textbf{линейности} оператора $T$:
\[ T(x) = T\left(\sum_{j=1}^n x_j v_j\right) = \sum_{j=1}^n x_j T(v_j). \]
Каждый вектор $T(v_j) \in W$ раскладывается по базису $\mathcal{C}$:
\[ T(v_j) = \sum_{i=1}^m a_{ij} w_i, \quad \text{то есть } [T(v_j)]_\mathcal{C} =
\begin{pmatrix}
            a_{1j} \\ a_{2j} \\ \vdots \\ a_{mj}
\end{pmatrix} \].
Подставим это выражение в формулу для $T(x)$:
\[ T(x) = \sum_{j=1}^n x_j \left( \sum_{i=1}^m a_{ij} w_i \right) = \sum_{i=1}^m \left( \sum_{j=1}^n a_{ij} x_j
\right) w_i.
\] Следовательно , $i$-я координата вектора $T(x)$ в базисе $\mathcal{C}$ равна $y_i = \sum_{j=1}^n a_{ij} x_j$.
В терминах матричного произведения это тождественно равно:
\[ \begin{pmatrix}
       y_1 \\ \vdots \\ y_m
\end{pmatrix} = \begin{pmatrix}
                    a_{11} & a_{12} & \dots & a_{1n} \\
                    a_{21} & a_{22} & \dots & a_{2n} \\
                    \vdots & \vdots & \ddots & \vdots \\
                    a_{m1} & a_{m2} & \dots & a_{mn}
\end{pmatrix} \begin{pmatrix}
                  x_1 \\ \vdots \\ x_n
\end{pmatrix} = A [x]_\mathcal{B}. \]

\textbf{Вывод:} Если мы возьмём в качестве $x$ $j$-й базисный вектор $v_j$, его координатный столбец $[v_j]_\mathcal{B}
= e_j = (0, \dots, 1, \dots, 0)^\top$ содержит единицу на $j$-м месте и нули на остальных.
Тогда:
\[ [T(v_j)]_\mathcal{C} = A e_j = A_{\cdot j} \quad (\text{$j$-й столбец матрицы } A). \]
Таким образом , $j$-й столбец матрицы отображения обязан в точности совпадать с вектором координат образа $j$-го
базисного вектора.

% -------------------------------------------------------------------------

\subsection{Теоретический вопрос II.2: Концептуально объясните rank-nullity}
\textbf{Формулировка вопроса:} Объясните концептуальный смысл теоремы о ранге и дефекте ($\dim V = \rank T + \nullity
T$) как закона сохранения размерности.

\textbf{Ответ:}

Теорема о ранге и дефекте формулируется следующим образом:
\[ \dim V = \dim(\ker T) + \dim(\im T) \iff \dim V = \nullity(T) + \rank(T). \]
Для матрицы $A \in \mathbb{R}^{m \times n}$ с $n$ столбцами:
\[ n = \dim(\Null A) + \rank(A). \]

\textbf{Концептуальная интерпретация («Закон сохранения степеней свободы»):}
\begin{enumerate}
    \item Пространство прообразов $V$ обладает $\dim V = n$ независимыми степенями свободы (направлениями движения).
    \item При действии любого линейного оператора $T$ эти степени свободы расщепляются на две строго
    непересекающиеся группы:
    \begin{itemize}
        \item Направления, которые отображение \textbf{уничтожает (схлопывает в нуль)}.
        Это ядро оператора $\ker T$.
        Количество схлопнутых измерений равно дефекту $\nullity(T)$.
        \item Направления, которые \textbf{сохраняются (переживают трансформацию)}.
        Они разворачиваются в образе $\im T$, формируя его размерность — ранг $\rank(T)$.
    \end{itemize}
    \item \textbf{Закон сохранения:} Измерение не может возникнуть из ниоткуда и не может исчезнуть бесследно:
    \[ \text{Исходные измерения} = \text{Потерянные измерения} + \text{Выжившие измерения}. \]
\end{enumerate}

\textbf{Экстремальные случаи:}
\begin{itemize}
    \item Если $\nullity(T) = 0$ (ядро тривиально, $\ker T = \{0\}$), то оператор инъективен: ни одно измерение
    не потеряно, $\rank(T) = n$.
    \item Если $\rank(T) = 0$ (нулевой оператор), все $n$ измерений схлопываются в точку $0$, $\nullity(T) = n$.
\end{itemize}

% -------------------------------------------------------------------------

\subsection{Теоретический вопрос II.3: Почему elementary row operations не искажают решения системы уравнений?}
\textbf{Формулировка вопроса:} Докажите, почему элементарные строчные преобразования строго сохраняют множество решений
системы линейных уравнений.

\textbf{Ответ и строгое доказательство:}

Рассмотрим систему линейных уравнений в матричной форме:
\[ Ax = b, \quad A \in \mathbb{R}^{m \times n}, \; x \in \mathbb{R}^n, \; b \in \mathbb{R}^m. \]
Обозначим множество ее решений через $S = \{ x \in \mathbb{R}^n \mid Ax = b \}$.
Элементарные преобразования строк включают: 1) перестановку строк $i$ и $j$; 2) умножение строки на $c \ne 0$; 3)
прибавление к строке $i$ строки $j$, умноженной на $c$.

\textbf{Лемма об элементарных матрицах:} Каждое элементарное строчное преобразование расширенной матрицы $(A \mid b)$ в
точности эквивалентно умножению ее слева на соответствующую \textbf{элементарную матрицу} $E \in \mathbb{R}^{m \times
m}$:
\[ E (A \mid b) = (EA \mid Eb). \]
Преобразованная система имеет вид:
\[ (EA) x = Eb. \]
Обозначим множество ее решений через $S' = \{ x \in \mathbb{R}^n \mid (EA)x = Eb \}$.

\textbf{Ключевое свойство:} Любая элементарная матрица $E$ \textbf{строго обратима}, причём матрица $E^{-1}$ сама
является элементарной матрицей того же типа:
\begin{itemize}
    \item Для перестановки строк $E_{ij}^{-1} = E_{ij}$;
    \item Для масштабирования $E_i(c)^{-1} = E_i(1/c)$ при $c \ne 0$;
    \item Для сдвига $E_{ij}(c)^{-1} = E_{ij}(-c)$.
\end{itemize}

\textbf{Доказательство эквивалентности множеств решений($S = S'$):}
\begin{enumerate}
    \item \textbf{Включение $S \subseteq S'$:} Пусть $x \in S$, то есть $Ax = b$.
    \item Умножим обе части равенства слева на матрицу $E$:
    \[ E(Ax) = Eb \iff (EA)x = Eb \implies x \in S'. \]
    \item \textbf{Включение $S' \subseteq S$:} Пусть $x \in S'$, то есть $(EA)x = Eb$.
    В силу существования $E^{-1}$ умножим обе части равенства слева на $E^{-1}$:
    \[ E^{-1} ((EA)x)=E^{-1} (Eb) \iff (E^{-1} E) Ax=(E^{-1} E) b \iff I_m Ax=I_m b \iff Ax = b \implies x \in S. \]
\end{enumerate}
Поскольку $S \subseteq S'$ и $S' \subseteq S$, множества решений тождественно совпадают: $S = S'$.
Ни одно решение не теряется и ни одно постороннее решение не появляется.

\newpage